\chapter{Worked example: adding a schema version}\label{ch:schema-version}

The Test Run record is defined in one file, \texttt{src/\allowbreak{}ate/\allowbreak{}storage/\allowbreak{}schema.py}, and one number:
\texttt{SCHEMA\_VERSION}. Everything that reads or writes the record goes through
\texttt{RunStore}, so a table added here is the only way a table exists. The example below is run by
\texttt{tests/\allowbreak{}test\_\allowbreak{}manuals.py}: it migrates a database with the shipped code,
tries a copy of the code that bumps only the number, migrates the same database with a copy that
declares version 6 and one more table, and then shows the shipped code refusing that database.

\section{The two edits}

Add the table to the metadata, and bump the version. The comment that explains each bump is part of
the convention: the file's own history is written in those comments --- \texttt{input\_value} made it
2, \texttt{media\_release} and \texttt{run\_file} made it 3, \texttt{schedule\_instance} made it 4,
\texttt{equipment} made it 5.

\begin{lstlisting}[language=ATEPython]
SCHEMA_VERSION = 6
\end{lstlisting}

\begin{lstlisting}[language=ATEPython]
review_note = Table(
    "review_note",
    metadata,
    Column("run_uuid", Text, primary_key=True),
    Column("seq", Integer, primary_key=True),  # the composite key, as every other run table
    Column("text", Text, nullable=True),
    Column("at", Text, nullable=True),
)
\end{lstlisting}

Nothing else is required: no hand-written DDL, no migration script. The version lives in
\texttt{PRAGMA user\_version}, and the tables are created by
\texttt{metadata.create\_\allowbreak{}all(engine)} in \texttt{RunStore.\_\allowbreak{}ensure\_\allowbreak{}schema}, which reflects each
table through the backend's own catalog and so never raises on a database that already has it.

The types are the ones the record uses everywhere: UUIDs as 36-character \texttt{Text}, times as
ISO-8601 \texttt{Text}, booleans as \texttt{Integer} 0/1, structured values as JSON \texttt{Text}.
There is deliberately no SQLite-specific column: a future Postgres backend is a connection change,
not a schema change. The only raw SQL in the storage layer is the two pragmas (WAL mode and
\texttt{user\_version}), and they stay in \texttt{store.py}.

\section{What a migration does, and what it does not}

\begin{lstlisting}[language=ATEShell]
python -m ate db migrate /tmp/station/data/ate.db --yes
\end{lstlisting}

prints \texttt{backed up \dots to \dots.backup-<YYYYMMDD-HHMMSS>} and then
\texttt{created or upgraded database at \dots (schema version 6)}. In the test this is run against the
same file more than once: with the shipped code it creates the ten tables the record has and makes no
copy, because there was no database to copy; with the patched copy it takes the version 5 database as
a copy beside itself, then adds \texttt{review\_note} on top. The assertion that matters is that the
set of tables only grew --- a migration adds, it never rewrites or drops --- and that the copy is the
database as it was, version 5 and no \texttt{review\_note}.

The last run is the reason the version check exists: the shipped code, pointed at the version 6
database, exits 1 with \texttt{this ATE is too old for \dots; update ATE}. \texttt{check\_\allowbreak{}database\_\allowbreak{}at\_\allowbreak{}launch}
compares \texttt{PRAGMA user\_version} with \texttt{SCHEMA\_VERSION} and returns
\texttt{NEEDS\_MIGRATE} when the database is behind, \texttt{TOO\_OLD} when it is ahead, and
\texttt{OK} only when they match. A run against a database whose tables are missing fails at the
first insert; the check is what makes the failure a message instead of an \texttt{OperationalError}
in a worker.

Spec rule 20 asked for a timestamped backup before an upgrade. It is no longer something the
Administrator has to do by hand: \texttt{ate db migrate} and \texttt{ate setup} both go through
\texttt{src/ate/storage/backup.py}, which copies the database beside itself before a migration runs.
The admin manual's database chapter says what an Administrator still has to back up themselves.

\section{Adding a column instead of a table}

\texttt{create\_\allowbreak{}all} does not alter a table that already exists, so adding a column to
\texttt{node} with a version bump leaves existing databases without it. Packet 35 closed the quiet
version of that mistake, the one where nothing at all is created: \texttt{RunStore.\_ensure\_\allowbreak{}schema}
stamps \texttt{PRAGMA user\_\allowbreak{}version} only when \texttt{missing\_\allowbreak{}tables} found something to
create. A database that stores an older version but already holds every table this ATE declares gets
no stamp and no silent pass. It has its own launch status, \texttt{NO\_\allowbreak{}MIGRATION}, and
\texttt{ate db status}, \texttt{ate db migrate} and the launch dialog refuse it with the same line:

\begin{verbatim}
the database at /var/lib/ate/ate.db stores schema version 5 and already holds
every table this ATE declares: this ATE has no migration from version 5 to 6,
so it is not stamped as version 6, and `ate db migrate` on it changes nothing
\end{verbatim}

The test runs exactly that: the shipped code, with only the number changed, exits 1 with that
sentence, and the database still says version 5 afterwards. Nothing in
\texttt{src/ate} writes an \texttt{ALTER} or an \texttt{add\_column}, and every insert names its columns
explicitly (\texttt{insert\_test\_run} builds its row field by field), so a new column has to be written
and tested as an explicit migration. The version number is a guard, not a migration.

What packet 35 did not close is the case the guard cannot see. A version bump that adds a column
\emph{alongside} a new table is still stamped silently: \texttt{missing\_tables} finds the new table, so
the stamp is raised, and the column on the table that already existed is never added. The database
then claims version 6 with an old \texttt{node} in it, and nothing notices. Before the next schema
version there has to be a migration registry --- a list of steps, each from version N to N+1, run in
order, each tested against a database at N --- so that a version number means "these steps ran"
instead of "some table was created".

\section{Where a new field reaches the report}

A table is not a report. The HTML writer in \texttt{src/\allowbreak{}ate/\allowbreak{}report/\allowbreak{}html.py} reads the record through
the store's query methods; a new table that nothing queries is invisible on the printed page. The
batching rules in chapter~\ref{ch:storage} apply to a new row type as well:
\texttt{\_\allowbreak{}rows\_\allowbreak{}per\_\allowbreak{}statement} keeps an insert under SQLite's variable limit, and
\texttt{\_log\_seq} is what orders a new row against the rest of the run.

\cando{add a table, bump the version, and show what a database at a different version does.}
